% ==========================================================
% BODY TEMPLATE — LaTeX (memoir class)
% This file is included by interior.tex via \input{body}
%
% COMMANDS AVAILABLE:
%   \chapteropener{1}{Chapter Title}       — ornamental chapter opener
%   Start the first paragraph normally; drop caps are not used.
%   \begin{recipebox}{Name} ... \end{recipebox}  — styled recipe block
%   \begin{accentbox}{Title} ... \end{accentbox} — styled callout/tip box
%   \yieldline{Yield: 24 • Prep: 15 min}  — centered yield line
%   \ornamentaldivider                    — three-diamond divider
%   \thinnerule                           — thin accent rule
%
% \chapteropener starts a fresh page; do not add another page break before it.
%
% TABLE RULES (A5 margin safety):
%   - Anywhere:        \begin{tabularx}{\linewidth}{...} with X columns.
%   - Inside a box:    ALWAYS \linewidth, never \textwidth.
%   - Plain tabular:   sum of p-widths + 2*ncols*6pt must stay <= \linewidth.
%   Good: \begin{tabularx}{\linewidth}{lXX} ... \end{tabularx}
%   Bad:  \begin{tabular}{p{0.3\linewidth}p{0.3\linewidth}p{0.3\linewidth}}
%         (0.9\linewidth of columns + 36pt of padding overflows A5)
% ==========================================================

\chapteropener{1}{Chapter Title}

The first paragraph of your chapter starts normally. This is where you
introduce the chapter's theme, set the tone, and draw the reader in.

Subsequent paragraphs are indented and justified automatically by the
memoir class. No special markup is needed for body text.

\section{Section Heading}

Level-2 headings are styled with the accent color and a larger font size.
They appear in the table of contents and in the running headers.

\begin{recipebox}{Recipe Name}

\textbf{Ingredients:}
\begin{itemize}
  \item 1 cup all-purpose flour
  \item \(\frac{1}{2}\) cup sugar
  \item 2 large eggs
  \item 1 teaspoon vanilla extract
\end{itemize}

\textbf{Instructions:}
\begin{enumerate}
  \item Preheat the oven to 350°F (175°C).
  \item In a large bowl, combine the dry ingredients.
  \item In a separate bowl, whisk together the wet ingredients.
  \item Fold the wet into the dry until just combined.
  \item Pour into a prepared pan and bake for 25--30 minutes.
\end{enumerate}

\end{recipebox}

\yieldline{Yield: 12 servings \textbullet{} Prep: 15 minutes \textbullet{} Cook: 30 minutes}

\begin{accentbox}{Key Concept}
Callout boxes highlight tips, key concepts, or comparisons. Tables inside
any box must use \texttt{\textbackslash linewidth}, never
\texttt{\textbackslash textwidth}:

\begin{tabularx}{\linewidth}{lXX}
\toprule
\textbf{Feature} & \textbf{Option A} & \textbf{Option B} \\
\midrule
Cost & Low & High \\
Speed & Fast & Slow \\
\bottomrule
\end{tabularx}
\end{accentbox}

\thinnerule

More paragraph text after the recipe. The recipe box provides a visual
break and makes it easy for readers to find individual recipes when
flipping through the book.

\clearpage

\chapteropener{2}{Next Chapter}

Beginning another chapter with the ornamental chapter opener creates a
consistent rhythm throughout the book.
